% Generated draft: requires mathematical and source review.
\documentclass[11pt]{article}
\usepackage[margin=25mm]{geometry}
\usepackage{fontspec}
\setmainfont{Latin Modern Roman}
\usepackage{amsmath,amssymb,mathtools,unicode-math}
\setmathfont{Latin Modern Math}
\usepackage{enumitem,needspace,xurl,hyperref}
\setlength{\parindent}{0pt}
\setlength{\parskip}{5pt}
\setlength{\emergencystretch}{3em}
\newcommand{\N}{\mathbb{N}}
\newcommand{\Z}{\mathbb{Z}}
\newcommand{\Q}{\mathbb{Q}}
\newcommand{\R}{\mathbb{R}}
\newcommand{\C}{\mathbb{C}}
\newcommand{\F}{\mathbb{F}}
\newcommand{\A}{\mathbb{A}}
\newcommand{\PP}{\mathbb{P}}
\newcommand{\E}{\mathbb{E}}
\newcommand{\eps}{\varepsilon}
\newcommand{\dd}{\,\mathrm{d}}
\begin{document}
\section*{UnsolvedMath ID: 800767}
\textbf{Source problem code:} catalogue entry
\subsection{Definitions and mathematical statement}
Let \(X\) be a compact K\"ahler manifold. A K\"ahler metric on \(X\) is a Hermitian Riemannian metric \(g\) whose associated real \(2\)-form \(\omega(u,v)=g(Ju,v)\) is closed, where \(J\) is the complex structure. Its holomorphic sectional curvature in a nonzero real tangent direction \(v\in T_xX\) is
\[
H_g(v)=K_g\bigl(\operatorname{span}_{\R}\{v,Jv\}\bigr),
\]
where \(K_g\) denotes Riemannian sectional curvature. Positive holomorphic sectional curvature means
\[
H_g(v)>0\qquad\text{for every }x\in X\text{ and every }0\ne v\in T_xX.
\]
The hypothesis is the existence of such a K\"ahler metric; it does not quantify over all K\"ahler metrics on \(X\).

A complex manifold is projective if it admits a holomorphic embedding into \(\PP^N_{\C}\) for some finite integer \(N\). A positive-dimensional smooth projective variety is rationally connected if there is a nonempty Zariski-open subset \(U\subseteq X\) such that every pair of distinct points of \(U\) lies in the image of a morphism \(\PP^1_{\C}\to X\). A point is regarded as rationally connected.

\textbf{Yau's conjecture, as reproduced in Matsumura's Conjecture 2.} For every compact K\"ahler manifold \(X\), if \(X\) admits a K\"ahler metric with positive holomorphic sectional curvature, then \(X\) is projective and rationally connected.

The supplied statement does not explicitly specify connectedness. The formulation above uses the customary connected interpretation of \(X\); if disconnected manifolds are allowed, the asserted conclusion must instead be interpreted for each connected component. This qualification requires source review.
\subsection{Short English statement}
Yau's conjecture, reproduced by Matsumura in 2022, asserts that a compact connected K\"ahler manifold admitting an everywhere strictly positive holomorphic sectional curvature metric is projective and rationally connected. Connectedness is implicit in this formulation and is not explicit in the supplied source. The registry labels the problem open; that label is supplied curation, not an independently verified status assessment.
\subsection{Sources}
\begin{itemize}
\item Supplied local UnsolvedMath registry record: permanent ID \(800767\), source problem code \texttt{catalogue entry}, title ``Positively Curved Projective Varieties --- Conjecture 2,'' AMR Open Problem Lists. Upstream revision: \texttt{d371ee917ac9647fbb06cd6fc529deec6f3dda82}.
\item Matsumura, \emph{Open Problems on Structure of Positively Curved Projective Varieties} (2022), Conjecture 2. Supplied source URL: \url{https://arxiv.org/abs/2211.09153}. The source record reports access on 2026-10-16 and extraction by \texttt{arxiv\_latex}; these are supplied provenance statements, not independently verified access claims.
\item The supplied statement attributes the conjecture to Yau and cites explicitly ``[Yau82, Problem 47].'' No full bibliographic entry for this reference is supplied. The registry's proposer field lists Matsumura (2022), while the statement identifies Yau as the originator.
\item The supplied background describes the conjecture as an analogy of Mok's result for holomorphic sectional curvature. No precise statement or bibliographic reference for Mok's result is supplied.
\end{itemize}
\end{document}
